% --- Archon physics formalization source begin ---
% archon:physics
% archon:covers PhyXMiniProblems/problem_phyx_mini_0856.lean
% archon:source-report reports/phyx_mini/problem_phyx_mini_0856.source.json

\chapter{Physics problem phyx\_mini\_0856}
\label{ch:PhyXMiniProblems_problem_phyx_mini_0856}

\paragraph{Problem source.}
\#\# Physical scenario

The figure shows the potential energy of an electric dipole. Consider a dipole that oscillates between \textbackslash{}( \textbackslash{}pm 60\textasciicircum{}\textbackslash{}circ \textbackslash{}).

\#\# Question

What is the dipole's kinetic energy when it is aligned with the electric field?

\#\# Answer choices

- A: A: 1.1\textbackslash{}mu\textbackslash{}mathrm\{J\}

- B: B: 1.2\textbackslash{}mu\textbackslash{}mathrm\{J\}

- C: C: 1.3\textbackslash{}mu\textbackslash{}mathrm\{J\}

- D: D: 1.0\textbackslash{}mu\textbackslash{}mathrm\{J\}

\#\# Figure caption (auxiliary; use the image as primary evidence)

The image shows a graph with a sinusoidal curve that represents a function of the form \textbackslash{}( U(\textbackslash{}phi) \textbackslash{}). The details are:

- **Axes**: 
  - The horizontal axis represents the angle \textbackslash{}(\textbackslash{}phi\textbackslash{}), marked in degrees (ranging from \textbackslash{}(-180\textasciicircum{}\textbackslash{}circ\textbackslash{}) to \textbackslash{}(180\textasciicircum{}\textbackslash{}circ\textbackslash{})).
  - The vertical axis represents the value \textbackslash{}( U \textbackslash{}) in microjoules (\textbackslash{}(\textbackslash{}mu J\textbackslash{})), ranging from \textbackslash{}(-2\textbackslash{}) to \textbackslash{}(2\textbackslash{}).
  
- **Curve**: 
  - The curve is sinusoidal, starting from the top at \textbackslash{}(-180\textasciicircum{}\textbackslash{}circ\textbackslash{}), descending through zero at \textbackslash{}(0\textasciicircum{}\textbackslash{}circ\textbackslash{}), reaching a minimum at around \textbackslash{}(-2\textbackslash{}) microjoules, then ascending back to zero at \textbackslash{}(180\textasciicircum{}\textbackslash{}circ\textbackslash{}).

- **Labels**:
  - The vertical axis is labeled \textbackslash{}( U\textbackslash{}, (\textbackslash{}mu J)\textbackslash{}).
  - The horizontal axis is labeled \textbackslash{}(\textbackslash{}phi\textbackslash{}).

The graph likely represents a potential energy function over a range of angles.

\#\# Dataset metadata

Electrostatics; Physical Model Grounding Reasoning, Multi-Formula Reasoning

\paragraph{Recorded answer/context.}
D: D: 1.0\textbackslash{}mu\textbackslash{}mathrm\{J\}

\paragraph{Figure/image path.}
/root/proposal\_for\_physic/science-mango/phyx\_mini\_run/phyx\_data/test\_image/856.png

\paragraph{Formalization target.}
create a compiling Lean file with sorry bodies at `PhyXMiniProblems/problem\_phyx\_mini\_0856.lean`. The Lean declarations must preserve the physical quantities, dimensions or dimensional roles, figure labels, governing-law hypotheses, and final relation expressed by this problem.
Use Mathlib/Physlib names found through LeanExplore where available. If a physics API is missing, introduce faithful local abstractions rather than scalar placeholder aliases.

\begin{theorem}[Physics formalization target]
\label{thm:physics:phyx_mini_0856:target}
\lean{PhyXMiniProblems.ProblemPhyXMini0856.kineticEnergyWhenAligned_eq_one_microjoule}
\uses{def:physics:phyx-mini-0856:phyxminiproblems-problemphyxmini0856-energyinmicrojoules, def:physics:phyx-mini-0856:phyxminiproblems-problemphyxmini0856-alignedangle, def:physics:phyx-mini-0856:phyxminiproblems-problemphyxmini0856-oscillatingelectricdipole, def:physics:phyx-mini-0856:phyxminiproblems-problemphyxmini0856-matchesprimaryfigure, def:physics:phyx-mini-0856:phyxminiproblems-problemphyxmini0856-oscillatesbetweensixtydegrees, def:physics:phyx-mini-0856:phyxminiproblems-problemphyxmini0856-satisfieselectricdipoleoscillationlaws, def:physics:phyx-mini-0856:phyxminiproblems-problemphyxmini0856-recordeddatasetanswer, def:physics:phyx-mini-0856:phyxminiproblems-problemphyxmini0856-answermatchesalignedkineticenergy}
For a dipole oscillating between \(\pm60^\circ\), conservation of mechanical
energy gives \(1\,\mu\mathrm J\) of kinetic energy at alignment, answer D.
\end{theorem}
\begin{proof}
The potential law is \(U(\phi)=-pE\cos\phi\).  The figure gives
\(U(0)=-2\,\mu\mathrm J\), hence \(pE=2\,\mu\mathrm J\).  Therefore
\(U(60^\circ)=-1\,\mu\mathrm J\).  At the turning point the kinetic energy is
zero, so the conserved total energy is \(-1\,\mu\mathrm J\).  At alignment,
\[
 K(0)=E_{\rm total}-U(0)=(-1)-(-2)=1\,\mu\mathrm J.
\]
The displayed-energy map assigns this value to recorded choice D.
\end{proof}

% --- Archon declaration topology begin ---
\section{Declaration topology}
\begin{definition}[electric Dipole Moment Dimension]
  \label{def:physics:phyx-mini-0856:phyxminiproblems-problemphyxmini0856-electricdipolemomentdimension}
  \lean{PhyXMiniProblems.ProblemPhyXMini0856.electricDipoleMomentDimension}
  The dimension C L of the magnitude of an electric dipole moment
\end{definition}
\begin{proof}
  This is the declared model definition of electric Dipole Moment Dimension.
\end{proof}
\begin{definition}[electric Field Strength Dimension]
  \label{def:physics:phyx-mini-0856:phyxminiproblems-problemphyxmini0856-electricfieldstrengthdimension}
  \lean{PhyXMiniProblems.ProblemPhyXMini0856.electricFieldStrengthDimension}
  The dimension M L T⁻² C⁻¹ of electric-field strength
\end{definition}
\begin{proof}
  This is the declared model definition of electric Field Strength Dimension.
\end{proof}
\begin{definition}[Electric Dipole Moment Magnitude]
  \label{def:physics:phyx-mini-0856:phyxminiproblems-problemphyxmini0856-electricdipolemomentmagnitude}
  \lean{PhyXMiniProblems.ProblemPhyXMini0856.ElectricDipoleMomentMagnitude}
  \uses{def:physics:phyx-mini-0856:phyxminiproblems-problemphyxmini0856-electricdipolemomentdimension}
  A nonnegative, unit-independent electric-dipole-moment magnitude
\end{definition}
\begin{proof}
  This is the declared model definition of Electric Dipole Moment Magnitude.
\end{proof}
\begin{definition}[Electric Field Strength Magnitude]
  \label{def:physics:phyx-mini-0856:phyxminiproblems-problemphyxmini0856-electricfieldstrengthmagnitude}
  \lean{PhyXMiniProblems.ProblemPhyXMini0856.ElectricFieldStrengthMagnitude}
  \uses{def:physics:phyx-mini-0856:phyxminiproblems-problemphyxmini0856-electricfieldstrengthdimension}
  A nonnegative, unit-independent electric-field-strength magnitude
\end{definition}
\begin{proof}
  This is the declared model definition of Electric Field Strength Magnitude.
\end{proof}
\begin{definition}[nonnegative SIReadout]
  \label{def:physics:phyx-mini-0856:phyxminiproblems-problemphyxmini0856-nonnegativesireadout}
  \lean{PhyXMiniProblems.ProblemPhyXMini0856.nonnegativeSIReadout}
  Coherent-SI readout of a nonnegative dimensionful quantity
\end{definition}
\begin{proof}
  This is the declared model definition of nonnegative SIReadout.
\end{proof}
\begin{definition}[dipole Moment In Coulomb Meters]
  \label{def:physics:phyx-mini-0856:phyxminiproblems-problemphyxmini0856-dipolemomentincoulombmeters}
  \lean{PhyXMiniProblems.ProblemPhyXMini0856.dipoleMomentInCoulombMeters}
  \uses{def:physics:phyx-mini-0856:phyxminiproblems-problemphyxmini0856-electricdipolemomentmagnitude, def:physics:phyx-mini-0856:phyxminiproblems-problemphyxmini0856-nonnegativesireadout}
  Electric-dipole-moment magnitude in coulomb-metres
\end{definition}
\begin{proof}
  This is the declared model definition of dipole Moment In Coulomb Meters.
\end{proof}
\begin{definition}[electric Field Strength In Newtons Per Coulomb]
  \label{def:physics:phyx-mini-0856:phyxminiproblems-problemphyxmini0856-electricfieldstrengthinnewtonspercoulomb}
  \lean{PhyXMiniProblems.ProblemPhyXMini0856.electricFieldStrengthInNewtonsPerCoulomb}
  \uses{def:physics:phyx-mini-0856:phyxminiproblems-problemphyxmini0856-electricfieldstrengthmagnitude, def:physics:phyx-mini-0856:phyxminiproblems-problemphyxmini0856-nonnegativesireadout}
  Electric-field-strength magnitude in newtons per coulomb
\end{definition}
\begin{proof}
  This is the declared model definition of electric Field Strength In Newtons Per Coulomb.
\end{proof}
\begin{definition}[joules]
  \label{def:physics:phyx-mini-0856:phyxminiproblems-problemphyxmini0856-joules}
  \lean{PhyXMiniProblems.ProblemPhyXMini0856.joules}
  Construct a signed physical energy from its coherent-SI joule readout
\end{definition}
\begin{proof}
  This is the declared model definition of joules.
\end{proof}
\begin{definition}[microjoules]
  \label{def:physics:phyx-mini-0856:phyxminiproblems-problemphyxmini0856-microjoules}
  \lean{PhyXMiniProblems.ProblemPhyXMini0856.microjoules}
  \uses{def:physics:phyx-mini-0856:phyxminiproblems-problemphyxmini0856-joules}
  Construct a signed physical energy from its microjoule readout
\end{definition}
\begin{proof}
  This is the declared model definition of microjoules.
\end{proof}
\begin{definition}[energy In Joules]
  \label{def:physics:phyx-mini-0856:phyxminiproblems-problemphyxmini0856-energyinjoules}
  \lean{PhyXMiniProblems.ProblemPhyXMini0856.energyInJoules}
  Read a signed physical energy in joules
\end{definition}
\begin{proof}
  This is the declared model definition of energy In Joules.
\end{proof}
\begin{definition}[energy In Microjoules]
  \label{def:physics:phyx-mini-0856:phyxminiproblems-problemphyxmini0856-energyinmicrojoules}
  \lean{PhyXMiniProblems.ProblemPhyXMini0856.energyInMicrojoules}
  \uses{def:physics:phyx-mini-0856:phyxminiproblems-problemphyxmini0856-energyinjoules}
  Read a signed physical energy in microjoules
\end{definition}
\begin{proof}
  This is the declared model definition of energy In Microjoules.
\end{proof}
\begin{definition}[degrees To Angle]
  \label{def:physics:phyx-mini-0856:phyxminiproblems-problemphyxmini0856-degreestoangle}
  \lean{PhyXMiniProblems.ProblemPhyXMini0856.degreesToAngle}
  Convert a real-valued degree readout to a geometric angle
\end{definition}
\begin{proof}
  This is the declared model definition of degrees To Angle.
\end{proof}
\begin{definition}[aligned Angle]
  \label{def:physics:phyx-mini-0856:phyxminiproblems-problemphyxmini0856-alignedangle}
  \lean{PhyXMiniProblems.ProblemPhyXMini0856.alignedAngle}
  Alignment of the dipole moment with the electric field, φ = 0
\end{definition}
\begin{proof}
  This is the declared model definition of aligned Angle.
\end{proof}
\begin{definition}[sixty Degrees]
  \label{def:physics:phyx-mini-0856:phyxminiproblems-problemphyxmini0856-sixtydegrees}
  \lean{PhyXMiniProblems.ProblemPhyXMini0856.sixtyDegrees}
  \uses{def:physics:phyx-mini-0856:phyxminiproblems-problemphyxmini0856-degreestoangle}
  The magnitude of either stated turning angle, 60°
\end{definition}
\begin{proof}
  This is the declared model definition of sixty Degrees.
\end{proof}
\begin{definition}[anti Aligned Angle]
  \label{def:physics:phyx-mini-0856:phyxminiproblems-problemphyxmini0856-antialignedangle}
  \lean{PhyXMiniProblems.ProblemPhyXMini0856.antiAlignedAngle}
  \uses{def:physics:phyx-mini-0856:phyxminiproblems-problemphyxmini0856-degreestoangle}
  Anti-alignment, represented by the 180° tick
\end{definition}
\begin{proof}
  This is the declared model definition of anti Aligned Angle.
\end{proof}
\begin{definition}[Figure Axis]
  \label{def:physics:phyx-mini-0856:phyxminiproblems-problemphyxmini0856-figureaxis}
  \lean{PhyXMiniProblems.ProblemPhyXMini0856.FigureAxis}
  The two axes visible in the supplied graph
\end{definition}
\begin{proof}
  This is the declared model definition of Figure Axis.
\end{proof}
\begin{definition}[Figure Axis Quantity]
  \label{def:physics:phyx-mini-0856:phyxminiproblems-problemphyxmini0856-figureaxisquantity}
  \lean{PhyXMiniProblems.ProblemPhyXMini0856.FigureAxisQuantity}
  The physical or geometric quantity assigned to a graph axis
\end{definition}
\begin{proof}
  This is the declared model definition of Figure Axis Quantity.
\end{proof}
\begin{definition}[Figure Axis Unit]
  \label{def:physics:phyx-mini-0856:phyxminiproblems-problemphyxmini0856-figureaxisunit}
  \lean{PhyXMiniProblems.ProblemPhyXMini0856.FigureAxisUnit}
  \uses{def:physics:phyx-mini-0856:phyxminiproblems-problemphyxmini0856-microjoules}
  The unit printed beside a graph axis
\end{definition}
\begin{proof}
  This is the declared model definition of Figure Axis Unit.
\end{proof}
\begin{definition}[Angle Tick]
  \label{def:physics:phyx-mini-0856:phyxminiproblems-problemphyxmini0856-angletick}
  \lean{PhyXMiniProblems.ProblemPhyXMini0856.AngleTick}
  The three labeled angular ticks in image 856.png
\end{definition}
\begin{proof}
  This is the declared model definition of Angle Tick.
\end{proof}
\begin{definition}[Potential Energy Tick]
  \label{def:physics:phyx-mini-0856:phyxminiproblems-problemphyxmini0856-potentialenergytick}
  \lean{PhyXMiniProblems.ProblemPhyXMini0856.PotentialEnergyTick}
  The three labeled energy ticks in image 856.png
\end{definition}
\begin{proof}
  This is the declared model definition of Potential Energy Tick.
\end{proof}
\begin{definition}[Trace Shape]
  \label{def:physics:phyx-mini-0856:phyxminiproblems-problemphyxmini0856-traceshape}
  \lean{PhyXMiniProblems.ProblemPhyXMini0856.TraceShape}
  Qualitative shape of the plotted potential-energy trace
\end{definition}
\begin{proof}
  This is the declared model definition of Trace Shape.
\end{proof}
\begin{definition}[Trace Color]
  \label{def:physics:phyx-mini-0856:phyxminiproblems-problemphyxmini0856-tracecolor}
  \lean{PhyXMiniProblems.ProblemPhyXMini0856.TraceColor}
  Color of the plotted trace in the primary raster
\end{definition}
\begin{proof}
  This is the declared model definition of Trace Color.
\end{proof}
\begin{definition}[Dipole Potential Energy Figure]
  \label{def:physics:phyx-mini-0856:phyxminiproblems-problemphyxmini0856-dipolepotentialenergyfigure}
  \lean{PhyXMiniProblems.ProblemPhyXMini0856.DipolePotentialEnergyFigure}
  \uses{def:physics:phyx-mini-0856:phyxminiproblems-problemphyxmini0856-figureaxis, def:physics:phyx-mini-0856:phyxminiproblems-problemphyxmini0856-figureaxisquantity, def:physics:phyx-mini-0856:phyxminiproblems-problemphyxmini0856-figureaxisunit, def:physics:phyx-mini-0856:phyxminiproblems-problemphyxmini0856-angletick, def:physics:phyx-mini-0856:phyxminiproblems-problemphyxmini0856-potentialenergytick, def:physics:phyx-mini-0856:phyxminiproblems-problemphyxmini0856-traceshape, def:physics:phyx-mini-0856:phyxminiproblems-problemphyxmini0856-tracecolor}
  Typed data represented by the potential-energy graph
\end{definition}
\begin{proof}
  This is the declared model definition of Dipole Potential Energy Figure.
\end{proof}
\begin{definition}[External Field Model]
  \label{def:physics:phyx-mini-0856:phyxminiproblems-problemphyxmini0856-externalfieldmodel}
  \lean{PhyXMiniProblems.ProblemPhyXMini0856.ExternalFieldModel}
  The external-field model relevant to the standard dipole energy law
\end{definition}
\begin{proof}
  This is the declared model definition of External Field Model.
\end{proof}
\begin{definition}[Oscillating Electric Dipole]
  \label{def:physics:phyx-mini-0856:phyxminiproblems-problemphyxmini0856-oscillatingelectricdipole}
  \lean{PhyXMiniProblems.ProblemPhyXMini0856.OscillatingElectricDipole}
  \uses{def:physics:phyx-mini-0856:phyxminiproblems-problemphyxmini0856-electricdipolemomentmagnitude, def:physics:phyx-mini-0856:phyxminiproblems-problemphyxmini0856-electricfieldstrengthmagnitude, def:physics:phyx-mini-0856:phyxminiproblems-problemphyxmini0856-dipolepotentialenergyfigure, def:physics:phyx-mini-0856:phyxminiproblems-problemphyxmini0856-externalfieldmodel}
  Independent physical observables for the oscillating dipole. The energy functions are not defined from one another; their relation is imposed by the governing laws below
\end{definition}
\begin{proof}
  This is the declared model definition of Oscillating Electric Dipole.
\end{proof}
\begin{definition}[Matches Primary Figure]
  \label{def:physics:phyx-mini-0856:phyxminiproblems-problemphyxmini0856-matchesprimaryfigure}
  \lean{PhyXMiniProblems.ProblemPhyXMini0856.MatchesPrimaryFigure}
  \uses{def:physics:phyx-mini-0856:phyxminiproblems-problemphyxmini0856-microjoules, def:physics:phyx-mini-0856:phyxminiproblems-problemphyxmini0856-energyinmicrojoules, def:physics:phyx-mini-0856:phyxminiproblems-problemphyxmini0856-alignedangle, def:physics:phyx-mini-0856:phyxminiproblems-problemphyxmini0856-antialignedangle, def:physics:phyx-mini-0856:phyxminiproblems-problemphyxmini0856-oscillatingelectricdipole}
  Literal axis labels, ticks, and calibrated potential-energy readouts from the primary raster. These are figure/data premises and do not state the requested kinetic energy at alignment
\end{definition}
\begin{proof}
  This is the declared model definition of Matches Primary Figure.
\end{proof}
\begin{definition}[Oscillates Between Sixty Degrees]
  \label{def:physics:phyx-mini-0856:phyxminiproblems-problemphyxmini0856-oscillatesbetweensixtydegrees}
  \lean{PhyXMiniProblems.ProblemPhyXMini0856.OscillatesBetweenSixtyDegrees}
  \uses{def:physics:phyx-mini-0856:phyxminiproblems-problemphyxmini0856-energyinjoules, def:physics:phyx-mini-0856:phyxminiproblems-problemphyxmini0856-alignedangle, def:physics:phyx-mini-0856:phyxminiproblems-problemphyxmini0856-sixtydegrees, def:physics:phyx-mini-0856:phyxminiproblems-problemphyxmini0856-oscillatingelectricdipole}
  The prose statement that the dipole oscillates between ±60°. The two endpoints are turning points, so their kinetic energy vanishes; alignment lies on the traversed arc. These facts are initial/trajectory data, not the current conclusion
\end{definition}
\begin{proof}
  This is the declared model definition of Oscillates Between Sixty Degrees.
\end{proof}
\begin{definition}[Satisfies Electric Dipole Oscillation Laws]
  \label{def:physics:phyx-mini-0856:phyxminiproblems-problemphyxmini0856-satisfieselectricdipoleoscillationlaws}
  \lean{PhyXMiniProblems.ProblemPhyXMini0856.SatisfiesElectricDipoleOscillationLaws}
  \uses{def:physics:phyx-mini-0856:phyxminiproblems-problemphyxmini0856-dipolemomentincoulombmeters, def:physics:phyx-mini-0856:phyxminiproblems-problemphyxmini0856-electricfieldstrengthinnewtonspercoulomb, def:physics:phyx-mini-0856:phyxminiproblems-problemphyxmini0856-energyinjoules, def:physics:phyx-mini-0856:phyxminiproblems-problemphyxmini0856-oscillatingelectricdipole}
  The governing electrostatic and mechanical laws. The potential-energy law is U(φ) = -p E cos φ; conservation and kinetic-energy nonnegativity are asserted on the configurations actually traversed by the oscillation. No field below mentions the requested value 1 μJ
\end{definition}
\begin{proof}
  This is the declared model definition of Satisfies Electric Dipole Oscillation Laws.
\end{proof}
\begin{definition}[Answer Choice]
  \label{def:physics:phyx-mini-0856:phyxminiproblems-problemphyxmini0856-answerchoice}
  \lean{PhyXMiniProblems.ProblemPhyXMini0856.AnswerChoice}
  Labels of the four kinetic-energy choices printed in the source
\end{definition}
\begin{proof}
  This is the declared model definition of Answer Choice.
\end{proof}
\begin{definition}[displayed Kinetic Energy In Microjoules]
  \label{def:physics:phyx-mini-0856:phyxminiproblems-problemphyxmini0856-displayedkineticenergyinmicrojoules}
  \lean{PhyXMiniProblems.ProblemPhyXMini0856.displayedKineticEnergyInMicrojoules}
  \uses{def:physics:phyx-mini-0856:phyxminiproblems-problemphyxmini0856-answerchoice}
  Kinetic-energy readout, in microjoules, printed beside each choice
\end{definition}
\begin{proof}
  This is the declared model definition of displayed Kinetic Energy In Microjoules.
\end{proof}
\begin{definition}[recorded Dataset Answer]
  \label{def:physics:phyx-mini-0856:phyxminiproblems-problemphyxmini0856-recordeddatasetanswer}
  \lean{PhyXMiniProblems.ProblemPhyXMini0856.recordedDatasetAnswer}
  \uses{def:physics:phyx-mini-0856:phyxminiproblems-problemphyxmini0856-answerchoice}
  The answer label recorded by the source dataset; it is not a premise
\end{definition}
\begin{proof}
  This is the declared model definition of recorded Dataset Answer.
\end{proof}
\begin{definition}[Answer Matches Aligned Kinetic Energy]
  \label{def:physics:phyx-mini-0856:phyxminiproblems-problemphyxmini0856-answermatchesalignedkineticenergy}
  \lean{PhyXMiniProblems.ProblemPhyXMini0856.AnswerMatchesAlignedKineticEnergy}
  \uses{def:physics:phyx-mini-0856:phyxminiproblems-problemphyxmini0856-energyinmicrojoules, def:physics:phyx-mini-0856:phyxminiproblems-problemphyxmini0856-alignedangle, def:physics:phyx-mini-0856:phyxminiproblems-problemphyxmini0856-oscillatingelectricdipole, def:physics:phyx-mini-0856:phyxminiproblems-problemphyxmini0856-answerchoice, def:physics:phyx-mini-0856:phyxminiproblems-problemphyxmini0856-displayedkineticenergyinmicrojoules}
  A displayed choice matches the independently modeled aligned energy
\end{definition}
\begin{proof}
  This is the declared model definition of Answer Matches Aligned Kinetic Energy.
\end{proof}
% --- Archon declaration topology end ---
% --- Archon physics formalization source end ---
